\documentclass[a4paper,12pt]{article}

% ===========================================================================
%  Лабораторная работа 3.6.1 «Спектральный анализ электрических сигналов»
%
%  Сборка:  cd 3.6.1/report && pdflatex main.tex  (или latexmk -pdf main.tex)
%  Таблицы, графики и числа в тексте читаются из ../data/processed/
%  (*.csv и fit_macros.tex), которые готовит ../scripts/process_data.py
%  по отсчётам со скриншотов (../scripts/digitize_screens.py).
%
% ===========================================================================

% Общая преамбула репозитория; если её нет (сборка вне репозитория) --
% локальная резервная копия.

% ============================================================
% Общий преамбул для отчётов по лабораторным работам (labs_lpr)
% Подключается из report/main.tex командой:
%   % ============================================================
% Общий преамбул для отчётов по лабораторным работам (labs_lpr)
% Подключается из report/main.tex командой:
%   % ============================================================
% Общий преамбул для отчётов по лабораторным работам (labs_lpr)
% Подключается из report/main.tex командой:
%   \input{../../common/latex_template/preamble.tex}
% ============================================================

% Поддержка русского языка и кодировки
\usepackage[T2A]{fontenc}
\usepackage{lmodern}
\usepackage[utf8]{inputenc}
\usepackage[russian]{babel}

% Математические пакеты
\usepackage{amsmath, amssymb}

% Настройка полей
\usepackage{geometry}
\geometry{top=2cm, bottom=2cm, left=2.5cm, right=1.5cm}

% Колонтитулы
\usepackage{fancyhdr}
\usepackage{graphicx}

% Современные таблицы
\usepackage{tabularray}

% Работа с csv файлами
\usepackage{csvsimple-l3}
\usepackage{tikz}

% Графики
\usepackage{pgfplots}
\pgfplotsset{compat=1.18}

% Ссылки
% \usepackage{hyperref} % гипперссылки, нужно оформление вместо красного квадрата
\usepackage[russian]{cleveref} % умные ссылки, нужно переоформление для формул, рисунков и тд (вместо ф-л что-то нужно придумать)

% ============================================================

% Поддержка русского языка и кодировки
\usepackage[T2A]{fontenc}
\usepackage{lmodern}
\usepackage[utf8]{inputenc}
\usepackage[russian]{babel}

% Математические пакеты
\usepackage{amsmath, amssymb}

% Настройка полей
\usepackage{geometry}
\geometry{top=2cm, bottom=2cm, left=2.5cm, right=1.5cm}

% Колонтитулы
\usepackage{fancyhdr}
\usepackage{graphicx}

% Современные таблицы
\usepackage{tabularray}

% Работа с csv файлами
\usepackage{csvsimple-l3}
\usepackage{tikz}

% Графики
\usepackage{pgfplots}
\pgfplotsset{compat=1.18}

% Ссылки
% \usepackage{hyperref} % гипперссылки, нужно оформление вместо красного квадрата
\usepackage[russian]{cleveref} % умные ссылки, нужно переоформление для формул, рисунков и тд (вместо ф-л что-то нужно придумать)

% ============================================================

% Поддержка русского языка и кодировки
\usepackage[T2A]{fontenc}
\usepackage{lmodern}
\usepackage[utf8]{inputenc}
\usepackage[russian]{babel}

% Математические пакеты
\usepackage{amsmath, amssymb}

% Настройка полей
\usepackage{geometry}
\geometry{top=2cm, bottom=2cm, left=2.5cm, right=1.5cm}

% Колонтитулы
\usepackage{fancyhdr}
\usepackage{graphicx}

% Современные таблицы
\usepackage{tabularray}

% Работа с csv файлами
\usepackage{csvsimple-l3}
\usepackage{tikz}

% Графики
\usepackage{pgfplots}
\pgfplotsset{compat=1.18}

% Ссылки
% \usepackage{hyperref} % гипперссылки, нужно оформление вместо красного квадрата
\usepackage[russian]{cleveref} % умные ссылки, нужно переоформление для формул, рисунков и тд (вместо ф-л что-то нужно придумать)


% Пакеты, нужные именно этой работе (повторная загрузка безопасна)
\usepackage{subcaption}
\usepackage{circuitikz}

% --- Служебные макросы ------------------------------------------------------
% Короткий пропуск (число/слово, которое надо вписать после измерений)
\newcommand{\qq}{\textbf{???}}
% Пропуск, требующий осмысленного текста
\newcommand{\fillme}{\textbf{?? НЕОБХОДИМО ЗАПОЛНИТЬ ??}}
% Имя файла/пути моноширинным шрифтом. \detokenize делает безопасными
% символы «_», «#», «%» и т.п. -- экранировать их (\_) не нужно.
\newcommand{\filename}[1]{\texttt{\detokenize{#1}}}
% Осциллограмма/спектр из report/img/. Если файла ещё нет -- рамка-заглушка.
% #1 -- имя файла, #2 -- доля ширины окружения (\linewidth: внутри subfigure
% это ширина подрисунка, а не всей страницы).
% Ширина заглушки уменьшена на рамку, чтобы она не выступала за поле.
\newcommand{\labimage}[2]{%
    \IfFileExists{img/#1}{%
        \includegraphics[width=#2\linewidth]{img/#1}%
    }{%
        \fbox{\parbox[c][0.6\dimexpr#2\linewidth\relax][c]{%
                \dimexpr#2\linewidth-2\fboxsep-2\fboxrule\relax}{%
            \centering\small\filename{img/#1}\\[2mm]\fillme}}%
    }%
}

% Настройка колонтитулов (как в остальных работах репозитория)
\pagestyle{fancy}
\fancyhf{}
\fancyhead[L]{%
    \IfFileExists{../../common/latex_template/miptlogo.png}{%
        \includegraphics[height=1cm]{../../common/latex_template/miptlogo.png}%
    }{%
        \textbf{[ЛОГОТИП МФТИ]}%
    }%
}
\fancyhead[R]{\textbf{Емельянов ДА, Байдин ФМ}}
\fancyfoot[C]{\thepage}
\renewcommand{\headrulewidth}{0.4pt}
\setlength{\headheight}{33pt}

% Числовые результаты обработки (генерирует scripts/process_data.py)
% Сгенерировано scripts/process_data.py -- вручную не править
\newcommand{\hDevMean}{-0{,}5}
\newcommand{\hDevMeanVal}{-0{,}5}
\newcommand{\hDevMeanNum}{-0.47131}
\newcommand{\hDevMin}{-0{,}7}
\newcommand{\hDevMinVal}{-0{,}7}
\newcommand{\hDevMinNum}{-0.720025}
\newcommand{\hDevMax}{-0{,}4}
\newcommand{\hDevMaxVal}{-0{,}4}
\newcommand{\hDevMaxNum}{-0.350306}
\newcommand{\aOne}{281{,}5}
\newcommand{\aOneVal}{281{,}5}
\newcommand{\aOneNum}{281.5}
\newcommand{\tauEnv}{50{,}7 \pm 0{,}2}
\newcommand{\tauEnvVal}{50{,}7}
\newcommand{\tauEnvNum}{50.7017}
\newcommand{\nuMeandA}{50{,}0 \pm 0{,}1}
\newcommand{\nuMeandAVal}{50{,}0}
\newcommand{\nuMeandANum}{49.9956}
\newcommand{\nuMeandB}{100{,}0 \pm 0{,}1}
\newcommand{\nuMeandBVal}{100{,}0}
\newcommand{\nuMeandBNum}{99.992}
\newcommand{\evenMax}{0{,}35}
\newcommand{\evenMaxVal}{0{,}35}
\newcommand{\evenMaxNum}{0.354}
\newcommand{\evenRel}{0{,}008}
\newcommand{\evenRelVal}{0{,}008}
\newcommand{\evenRelNum}{0.00787437}
\newcommand{\nRatioMinA}{0{,}83}
\newcommand{\nRatioMinAVal}{0{,}83}
\newcommand{\nRatioMinANum}{0.827}
\newcommand{\nRatioThreeA}{0{,}992}
\newcommand{\nRatioThreeAVal}{0{,}992}
\newcommand{\nRatioThreeANum}{0.992}
\newcommand{\nRatioFiveA}{0{,}974}
\newcommand{\nRatioFiveAVal}{0{,}974}
\newcommand{\nRatioFiveANum}{0.974}
\newcommand{\screenEdge}{780}
\newcommand{\screenEdgeVal}{780}
\newcommand{\screenEdgeNum}{779.681}
\newcommand{\fsamp}{1562{,}3 \pm 0{,}9}
\newcommand{\fsampVal}{1562{,}3}
\newcommand{\fsampNum}{1562.3}
\newcommand{\fnyq}{781{,}2 \pm 0{,}4}
\newcommand{\fnyqVal}{781{,}2}
\newcommand{\fnyqNum}{781.152}
\newcommand{\kDnu}{1{,}000 \pm 0{,}009}
\newcommand{\kDnuVal}{1{,}000}
\newcommand{\kDnuNum}{0.999645}
\newcommand{\kDnuFree}{1{,}000 \pm 0{,}002}
\newcommand{\kDnuFreeVal}{1{,}000}
\newcommand{\kDnuFreeNum}{1.00043}
\newcommand{\bDnu}{-0{,}02 \pm 0{,}05}
\newcommand{\bDnuVal}{-0{,}02}
\newcommand{\bDnuNum}{-0.0233666}
\newcommand{\prodMin}{0{,}98}
\newcommand{\prodMinVal}{0{,}98}
\newcommand{\prodMinNum}{0.98}
\newcommand{\prodMax}{1{,}02}
\newcommand{\prodMaxVal}{1{,}02}
\newcommand{\prodMaxNum}{1.02}
\newcommand{\prodMean}{0{,}999 \pm 0{,}005}
\newcommand{\prodMeanVal}{0{,}999}
\newcommand{\prodMeanNum}{0.998571}
\newcommand{\kdnu}{0{,}9999 \pm 0{,}0003}
\newcommand{\kdnuVal}{0{,}9999}
\newcommand{\kdnuNum}{0.999878}
\newcommand{\kdnuFree}{1{,}0005 \pm 0{,}0006}
\newcommand{\kdnuFreeVal}{1{,}0005}
\newcommand{\kdnuFreeNum}{1.00055}
\newcommand{\bdnu}{-0{,}0004 \pm 0{,}0003}
\newcommand{\bdnuVal}{-0{,}0004}
\newcommand{\bdnuNum}{-0.000409459}
\newcommand{\kBurst}{1{,}00 \pm 0{,}03}
\newcommand{\kBurstVal}{1{,}00}
\newcommand{\kBurstNum}{0.995415}
\newcommand{\kBurstSpacing}{1{,}000 \pm 0{,}001}
\newcommand{\kBurstSpacingVal}{1{,}000}
\newcommand{\kBurstSpacingNum}{1}
\newcommand{\burstProd}{0{,}99 \pm 0{,}01}
\newcommand{\burstProdVal}{0{,}99}
\newcommand{\burstProdNum}{0.99}
\newcommand{\amRatio}{0{,}503 \pm 0{,}002}
\newcommand{\amRatioVal}{0{,}503}
\newcommand{\amRatioNum}{0.503178}
\newcommand{\amM}{1{,}006 \pm 0{,}004}
\newcommand{\amMVal}{1{,}006}
\newcommand{\amMNum}{1.00636}
\newcommand{\amNumod}{2{,}00 \pm 0{,}03}
\newcommand{\amNumodVal}{2{,}00}
\newcommand{\amNumodNum}{2.002}
\newcommand{\amNuzero}{50{,}02 \pm 0{,}05}
\newcommand{\amNuzeroVal}{50{,}02}
\newcommand{\amNuzeroNum}{50.023}
\newcommand{\amMain}{707{,}9 \pm 1{,}5}
\newcommand{\amMainVal}{707{,}9}
\newcommand{\amMainNum}{707.9}
\newcommand{\amSide}{356{,}2 \pm 1{,}5}
\newcommand{\amSideVal}{356{,}2}
\newcommand{\amSideNum}{356.2}
\newcommand{\filtT}{2{,}999 \pm 0{,}002}
\newcommand{\filtTVal}{2{,}999}
\newcommand{\filtTNum}{2.9992}
\newcommand{\filtTau}{0{,}145 \pm 0{,}005}
\newcommand{\filtTauVal}{0{,}145}
\newcommand{\filtTauNum}{0.145235}
\newcommand{\filtRC}{3{,}48 \pm 0{,}01}
\newcommand{\filtRCVal}{3{,}48}
\newcommand{\filtRCNum}{3.48376}
\newcommand{\filtUzero}{3{,}98 \pm 0{,}06}
\newcommand{\filtUzeroVal}{3{,}98}
\newcommand{\filtUzeroNum}{3.97753}
\newcommand{\filtNuRC}{45{,}7 \pm 0{,}1}
\newcommand{\filtNuRCVal}{45{,}7}
\newcommand{\filtNuRCNum}{45.6848}
\newcommand{\filtNuRCinv}{287 \pm 1}
\newcommand{\filtNuRCinvVal}{287}
\newcommand{\filtNuRCinvNum}{287.046}
\newcommand{\filtNuOne}{333{,}5 \pm 0{,}1}
\newcommand{\filtNuOneVal}{333{,}5}
\newcommand{\filtNuOneNum}{333.531}
\newcommand{\filtTspec}{2{,}998 \pm 0{,}001}
\newcommand{\filtTspecVal}{2{,}998}
\newcommand{\filtTspecNum}{2.99822}
\newcommand{\filtDevMedian}{17}
\newcommand{\filtDevMedianVal}{17}
\newcommand{\filtDevMedianNum}{17.4667}
\newcommand{\filtDevMax}{35}
\newcommand{\filtDevMaxVal}{35}
\newcommand{\filtDevMaxNum}{35.4485}
\newcommand{\filtWnu}{7{,}3}
\newcommand{\filtWnuVal}{7{,}3}
\newcommand{\filtWnuNum}{7.30069}
\newcommand{\filtQ}{0{,}048}
\newcommand{\filtQVal}{0{,}048}
\newcommand{\filtQNum}{0.0484244}


% Числа из CSV печатаются с десятичной запятой
\pgfkeys{/pgf/number format/.cd, use comma, 1000 sep={}, precision=4}
\newcommand{\pn}[1]{\pgfmathprintnumber{#1}}
\setlength{\emergencystretch}{2em}

\begin{document}

% --- ТИТУЛЬНЫЙ ЛИСТ ---
\begin{titlepage}
\begin{center}
    {\large МОСКОВСКИЙ ФИЗИКО-ТЕХНИЧЕСКИЙ ИНСТИТУТ}
\end{center}
\begin{center}
    {\large Физтех-школа физики и исследований им. Ландау\break ЛФИ}
\end{center}

    \vspace{3cm}

\vspace{0.15cm}
{
\begin{center}
    \huge \bf Лабораторная работа №3.6.1 по курсу Общая физика\\
    \vspace{0.4cm}
    \LARGE Спектральный анализ электрических сигналов.
\end{center}
}
\vspace{1.5cm}
\begin{flushright}
Емельянов Дмитрий Алексеевич\\
Байдин Фаддей Михайлович\\
\vspace{0.4cm}
Б02-502
\end{flushright}
\vspace{10cm}
\begin{center}
    \large 15 сентября 2026, г. Долгопрудный
\end{center}
\end{titlepage}
% ----------------------

\section{Аннотация}
В работе изучаются спектры периодических электрических сигналов различной
формы и влияние параметров сигнала на вид его спектра. С помощью генератора
сигналов формируются последовательности прямоугольных импульсов (в том числе
меандр), цугов синусоиды, а также синусоидальный сигнал, модулированный по
амплитуде; их спектры наблюдаются с помощью USB-осциллографа в режиме
быстрого преобразования Фурье (БПФ).

Измерены частоты и амплитуды гармоник, результаты
сравнены с рассчитанными по теории Фурье, проверены соотношения
неопределённостей $\Delta\nu\cdot\tau\simeq1$ и $\delta\nu\cdot T=1$,
теорема смещения и соотношение $a_{\text{бок}}/a_{\text{осн}} = m/2$ для
АМ-сигнала. На спектре меандра в широком диапазоне частот наблюдалось
наложение частот (теорема Котельникова) и определена частота дискретизации
спектроанализатора. По осциллограмме сигнала на выходе $RC$-цепочки
определена её постоянная времени $\tau_{RC}$, спектр выходного сигнала
сопоставлен с частотной характеристикой фильтра.

\section{Теоретические сведения}

\subsection{Ряд Фурье и спектр сигнала}
Согласно теореме Фурье, любая периодическая функция с периодом $T$
представима в виде ряда гармонических функций с кратными частотами:
\begin{equation}
    \label{eq:fourier}
    f(t) = \frac{A_0}{2}
    + \sum_{n=1}^{\infty}\bigl[A_n\cos(2\pi\nu_n t) + B_n\sin(2\pi\nu_n t)\bigr],
    \qquad \nu_n = \frac{n}{T},
\end{equation}
где коэффициенты разложения находятся как
\begin{equation}
    \label{eq:fourier_coef}
    A_n = \frac{2}{T}\int\limits_0^T f(t)\cos(2\pi\nu_n t)\,dt,
    \qquad
    B_n = \frac{2}{T}\int\limits_0^T f(t)\sin(2\pi\nu_n t)\,dt.
\end{equation}
Эквивалентная запись «в представлении амплитуд и фаз»:
\begin{equation}
    \label{eq:fourier_amp}
    f(t) = \frac{a_0}{2} + \sum_{n=1}^{\infty} a_n\cos(2\pi\nu_n t + \varphi_n),
    \qquad
    a_n = \sqrt{A_n^2 + B_n^2},\quad \operatorname{tg}\varphi_n = \frac{B_n}{A_n}.
\end{equation}
Совокупность частот $\nu_n$ и соответствующих им амплитуд $a_n$ называют
\emph{спектром} сигнала. Спектр периодического сигнала дискретен: гармоники
расположены на расстоянии $\delta\nu = 1/T$ друг от друга.

При измерениях на осциллографе доступны только амплитуды $|a_n|$ (информация
о фазах теряется), а единицы измерения амплитуд произвольны, поэтому
сравнивать теорию и эксперимент следует по \emph{относительным} величинам,
например $|a_n/a_1|$.

\subsection{Спектр последовательности прямоугольных импульсов}
Для периодической последовательности прямоугольных импульсов амплитуды
$V_0$, длительности $\tau$ и периода $T$ коэффициенты разложения равны
\begin{equation}
    \label{eq:rect_spectrum}
    \nu_n = \frac{n}{T},
    \qquad
    |a_n| = 2V_0\frac{\tau}{T}\,\frac{|\sin \pi\nu_n\tau|}{\pi\nu_n\tau}
    = \frac{2V_0}{\pi}\cdot\frac{\bigl|\sin\frac{\pi n\tau}{T}\bigr|}{n}.
\end{equation}
Огибающая спектра имеет вид $|\operatorname{sinc}(\pi\nu\tau)|$ и обращается в
ноль на частотах $\nu = k/\tau$, $k=1,2,\dots$ Относительные амплитуды,
удобные для сравнения с экспериментом:
\begin{equation}
    \label{eq:rect_ratio}
    \left|\frac{a_n}{a_1}\right|
    = \frac{\bigl|\sin\frac{\pi n \tau}{T}\bigr|}
           {n\,\bigl|\sin\frac{\pi\tau}{T}\bigr|}.
\end{equation}
Частный случай $\tau = T/2$ --- \emph{меандр}: все чётные гармоники
обращаются в ноль, а нечётные убывают как $1/n$:
\begin{equation}
    \label{eq:meander}
    |a_n| = \frac{2V_0}{\pi n}\quad (n = 1, 3, 5, \dots),
    \qquad
    \left|\frac{a_n}{a_1}\right| = \frac1n.
\end{equation}

\subsection{Соотношения неопределённостей}
Если у сигнала есть характерное время $\Delta t$, то в спектре появляется
характерный масштаб $\Delta\nu\sim 1/\Delta t$:
\begin{equation}
    \label{eq:uncertainty}
    \Delta\nu\cdot\Delta t \sim 1.
\end{equation}
Для последовательности прямоугольных импульсов это соотношение принимает два
конкретных вида:
\begin{equation}
    \label{eq:uncertainty_rect}
    \Delta\nu = \frac{1}{\tau},
    \qquad
    \delta\nu = \frac{1}{T},
\end{equation}
где $\Delta\nu$ --- ширина спектра (от $\nu=0$ до первого нуля огибающей),
а $\delta\nu$ --- расстояние между соседними гармониками. Второе соотношение
точное и от формы сигнала не зависит, первое --- по порядку величины и
зависит от того, что именно принято за «ширину».

\subsection{Теорема смещения. Спектр последовательности цугов}
Умножение сигнала $g(t)$ на гармоническую несущую $\cos 2\pi\nu_0 t$ приводит
к смещению его спектра: спектр произведения представляет собой спектр
огибающей $g(t)$, сдвинутый к частоте $\nu_0$ (и $-\nu_0$) и уменьшенный вдвое
по амплитуде:
\begin{equation}
    \label{eq:shift}
    g(t)\cos 2\pi\nu_0 t
    \ \longleftrightarrow\
    \tfrac12\, a_g(\nu-\nu_0) + \tfrac12\, a_g(\nu+\nu_0).
\end{equation}
Последовательность цугов --- это произведение синусоиды с частотой
$\nu_0$ и амплитудой $A$ на последовательность прямоугольных импульсов
длительности $\tau = N/\nu_0$ ($N$ --- число периодов синусоиды в цуге) с
периодом повторения $T$. Поэтому её спектр --- «сдвинутый» спектр
прямоугольных импульсов: гармоники расположены с шагом $\delta\nu = 1/T$,
огибающая $A\frac{\tau}{T}\,|\operatorname{sinc}\pi(\nu-\nu_0)\tau|$ имеет
максимум вблизи $\nu_0$, а расстояние от центра спектра до первого нуля
огибающей равно
\begin{equation}
    \label{eq:burst_width}
    \Delta\nu = \frac{1}{\tau} = \frac{\nu_0}{N}.
\end{equation}

\subsection{Амплитудная модуляция}
Амплитудно-модулированный сигнал с глубиной модуляции $m$ раскладывается на
три гармоники:
\begin{equation}
    \label{eq:am}
    f(t) = A_0\bigl(1 + m\cos 2\pi\nu_{\text{мод}}t\bigr)\cos 2\pi\nu_0 t
    = A_0\cos 2\pi\nu_0 t
    + \frac{mA_0}{2}\sum_{\pm}\cos 2\pi(\nu_0\pm\nu_{\text{мод}})t,
\end{equation}
откуда отношение амплитуд боковой и основной линий
\begin{equation}
    \label{eq:am_ratio}
    \frac{a_{\text{бок}}}{a_{\text{осн}}} = \frac{m}{2},
    \qquad\text{причём}\qquad
    m = \frac{A_{\max}-A_{\min}}{A_{\max}+A_{\min}}.
\end{equation}

\subsection{Дискретизация сигнала и наложение частот}
Цифровой спектроанализатор оцифровывает сигнал с частотой
$\nu_{\text{дискр}}$ и надёжно воспроизводит лишь составляющие с
$\nu < \nu_{\text{дискр}}/2$ (теорема Котельникова). Составляющая с
частотой $\nu > \nu_{\text{дискр}}/2$ не пропадает, а после дискретизации
неотличима от составляющей с частотой
\begin{equation}
    \label{eq:alias}
    \nu' = \bigl|\nu - k\,\nu_{\text{дискр}}\bigr|,
    \qquad k = 1, 2, \dots,
    \qquad 0 \leqslant \nu' \leqslant \nu_{\text{дискр}}/2,
\end{equation}
т.\,е. «зеркально отражается» в видимую полосу (наложение, англ.
\emph{aliasing}).

\subsection{Фильтрация сигнала $RC$-цепочкой}
$RC$-цепочка (\cref{fig:rc_scheme}) --- фильтр низких частот. Для
гармонического сигнала частоты $\nu$ отношение амплитуд напряжения на
конденсаторе и на входе равно
\begin{equation}
    \label{eq:rc_filter}
    K(\nu) = \left|\frac{1/i\omega C}{R + 1/i\omega C}\right|
    = \frac{1}{\sqrt{1 + (2\pi\nu\tau_{RC})^2}},
    \qquad \tau_{RC} = RC.
\end{equation}
Согласно принципу суперпозиции, каждая гармоника сигнала фильтруется
независимо: высокочастотные гармоники подавляются сильнее низкочастотных,
поэтому резкие фронты импульсов сглаживаются. При $2\pi\nu\tau_{RC}\gg1$
$K(\nu)\approx 1/(2\pi\nu\tau_{RC})$ --- цепочка работает как
интегратор, и амплитуды гармоник выходного сигнала убывают как $1/n$
быстрее, чем у входного.

Во временном представлении на каждом участке, где входное напряжение
постоянно и равно $U_\infty$, напряжение на конденсаторе релаксирует к нему
экспоненциально:
\begin{equation}
    \label{eq:rc_transient}
    U(t) = U_\infty + \bigl(U(0) - U_\infty\bigr)\,e^{-t/\tau_{RC}}.
\end{equation}

\section{Используемое оборудование}
\begin{itemize}
    \item генератор сигналов произвольной формы АКИП-3409/4;
    \item функциональный генератор WaveStation 2012;
    \item цифровой осциллограф SIGLENT АКИП-4131/1 с функцией БПФ;
    \item USB-осциллограф АКИП-4107 с программой PicoScope-6 (режимы
          осциллографа и спектроанализатора);
    \item $RC$-цепочка (фильтр низких частот);
    \item соединительные кабели.
\end{itemize}

\section{Методика измерений}
Выход генератора соединяется с входным каналом осциллографа. Сигнал
настраивается по временной развёртке (устойчивая картина, сигнал целиком
помещается на экране), после чего включается режим спектроанализатора.

Существенно, что спектр вычисляется только по той части сигнала, которая
попала на экран. Поэтому масштаб по вертикали подбирается так, чтобы сигнал
не выходил за пределы экрана (иначе он будет «обрезан», и спектр исказится),
а развёртка выбирается так, чтобы на экране умещалось много периодов
сигнала. Частота дискретизации ограничивает сверху доступный диапазон
частот: согласно теореме Котельникова надёжно измеряется область
$\nu < \nu_{\text{дискр}}/2$.

Частоты $\nu_n$ и амплитуды $|a_n|$ спектральных линий, ширина спектра
$\Delta\nu$ и расстояние между соседними линиями $\delta\nu$ в частях А и Б
измерялись с помощью курсоров спектроанализатора. Ширина спектра
определялась по номеру гармоники с нулевой амплитудой, поэтому её
погрешность принималась равной половине шага гармоник $\delta\nu/2$.
Количественные данные частей А и Б (\cref{tab:harmonics,tab:width,tab:spacing,tab:bursts})
получены на той же установке другой бригадой и используются с согласия
автора и преподавателя; спектрограммы на рисунках --- собственные.

Для спектров меандра, АМ-сигнала и сигнала на выходе $RC$-цепочки
частоты и амплитуды снимались с сохранённых изображений экрана
спектроанализатора по шкалам осей (скрипт
\filename{scripts/digitize_screens.py}, отсчёты хранятся в
\filename{data/raw/screens/}). Погрешность такого отсчёта принималась
равной цене одного пикселя изображения: по частоте $0{,}75$~кГц (спектр
в диапазоне $0\div780$~кГц), $0{,}05$~кГц (АМ), $3$~кГц ($RC$-цепочка); по
амплитуде $0{,}09$~мВ, $1{,}5$~мВ и $0{,}1$~мВ соответственно.

Амплитуды приведены в единицах шкалы осциллографа (мВ); для проверки
теории используются их отношения. Параметры аппроксимаций и их
погрешности находились методом наименьших квадратов (скрипт
\filename{scripts/process_data.py}).

В части Е собирается схема, показанная на \cref{fig:rc_scheme}: сигнал с
генератора подаётся на резистор $R$, а напряжение с конденсатора $C$ --- на
вход осциллографа.

\begin{figure}[!ht]
    \centering
    \begin{circuitikz}[scale=1.0, transform shape]
        \draw (0,0) to[sV, l=$U_{\text{вх}}$] (0,2.6);
        \draw (0,2.6) -- (1,2.6) to[R, l=$R$] (3.5,2.6) -- (5,2.6);
        \draw (3.5,2.6) to[C, l=$C$] (3.5,0);
        \draw (0,0) -- (5,0);
        \draw (5,2.6) to[open, v^>=$U_{\text{вых}}$] (5,0);
        \node at (0,-0.5) {\small генератор};
        \node at (5.2,-0.5) {\small осциллограф};
    \end{circuitikz}
    \caption{Схема установки для изучения фильтрации сигналов
             ($RC$-фильтр низких частот)}
    \label{fig:rc_scheme}
\end{figure}

\section{Результаты измерений и обработка данных}

\subsection{Параметры сигналов}
\begin{itemize}
    \item амплитуды гармоник: $T = 1$~мс ($\nu_{\text{повт}} = 1$~кГц),
          $\tau = T/20 = 50$~мкс;
    \item измерение $\Delta\nu(1/\tau)$: $T = 1$~мс,
          $\tau = 20\div200$~мкс ($T/50\div T/5$);
    \item измерение $\delta\nu(1/T)$: $\tau = 100$~мкс,
          $T = 2\tau\div50\tau$ ($0{,}2\div5$~мс);
    \item спектрограммы прямоугольных импульсов: $\tau = 50$~мкс,
          $\nu_{\text{повт}} = 1;\ 2;\ 3;\ 5$~кГц и $\nu_{\text{повт}} = 2$~кГц,
          $\tau = 50\div250$~мкс;
    \item меандр ($\tau = T/2$) с частотой $50$ и $100$~кГц,
          диапазон спектроанализатора $0\div780$~кГц;
    \item цуги: $\nu_0 = 30\div50$~кГц, $N = 5\div20$,
          $T = 1\div4$~мс;
    \item АМ-сигнал: $\nu_0 = 50$~кГц, $\nu_{\text{мод}} = 2$~кГц;
    \item $RC$-цепочка: прямоугольные импульсы с $T \approx 3$~мкс,
          $\tau \approx T/20$.
\end{itemize}

\subsection{А. Качественное исследование спектра прямоугольных импульсов}
Спектры последовательности прямоугольных импульсов при различных параметрах
показаны на \cref{fig:rect_spectra_nu,fig:rect_spectra_tau}.
\begin{itemize}
    \item При увеличении частоты повторения $\nu_{\text{повт}}$ и
          фиксированной длительности $\tau = 50$~мкс
          (\cref{fig:rect_spectra_nu}) расстояние между компонентами спектра
          растёт: линии стоят на частотах, кратных $1$, $2$, $3$ и $5$~кГц,
          т.\,е. $\delta\nu = \nu_{\text{повт}} = 1/T$. Одновременно растут
          амплитуды гармоник: при фиксированном $\tau$ они пропорциональны
          скважности $\tau/T$ (\cref{eq:rect_spectrum}). При
          $\nu_{\text{повт}} \geqslant 3$~кГц первая гармоника выходит за
          верхнюю границу экрана ($20{,}1$~мВ), а при
          $\nu_{\text{повт}} = 1$~кГц спектр снят в более чувствительном
          масштабе ($4{,}7$~мВ на весь экран) и зашкаливают все линии.
    \item При увеличении длительности импульса $\tau$ и фиксированной
          $\nu_{\text{повт}} = 2$~кГц (\cref{fig:rect_spectra_tau})
          расстояние между гармониками не меняется, а огибающая спектра
          сужается: амплитуды спадают с номером гармоники тем быстрее, чем
          длиннее импульс. При $\tau = 250$~мкс $= T/2$ первый нуль
          огибающей $1/\tau = 4$~кГц совпадает со второй гармоникой, и она
          исчезает из спектра.
    \item Между основными линиями видны слабые побочные пики (не более
          $1{,}7$~мВ, т.\,е. $\lesssim 10$~\% от амплитуды гармоник),
          положение которых не зависит от параметров сигнала; это
          паразитные составляющие («краевые эффекты» окна БПФ и наводки),
          в обработке они не учитывались.
\end{itemize}

\begin{figure}[!ht]
    \centering
    \begin{subfigure}{0.48\textwidth}
        \centering
        \labimage{spectrum_rect_nu1.png}{1.0}
        \caption{$\nu_{\text{повт}} = 1$~кГц (шкала $4{,}7$~мВ)}
    \end{subfigure}
    \hfill
    \begin{subfigure}{0.48\textwidth}
        \centering
        \labimage{spectrum_rect_nu2.png}{1.0}
        \caption{$\nu_{\text{повт}} = 2$~кГц}
    \end{subfigure}

    \vspace{4mm}
    \begin{subfigure}{0.48\textwidth}
        \centering
        \labimage{spectrum_rect_nu3.png}{1.0}
        \caption{$\nu_{\text{повт}} = 3$~кГц}
    \end{subfigure}
    \hfill
    \begin{subfigure}{0.48\textwidth}
        \centering
        \labimage{spectrum_rect_nu4.png}{1.0}
        \caption{$\nu_{\text{повт}} = 5$~кГц}
    \end{subfigure}
    \caption{Спектры последовательности прямоугольных импульсов при
             $\tau = 50$~мкс и различных частотах повторения}
    \label{fig:rect_spectra_nu}
\end{figure}

\begin{figure}[!ht]
    \centering
    \begin{subfigure}{0.48\textwidth}
        \centering
        \labimage{spectrum_rect_nu2.png}{1.0}
        \caption{$\tau = 50$~мкс}
    \end{subfigure}
    \hfill
    \begin{subfigure}{0.48\textwidth}
        \centering
        \labimage{spectrum_rect_tau100.png}{1.0}
        \caption{$\tau = 100$~мкс}
    \end{subfigure}

    \vspace{4mm}
    \begin{subfigure}{0.48\textwidth}
        \centering
        \labimage{spectrum_rect_tau200.png}{1.0}
        \caption{$\tau = 200$~мкс}
    \end{subfigure}
    \hfill
    \begin{subfigure}{0.48\textwidth}
        \centering
        \labimage{spectrum_rect_tau250.png}{1.0}
        \caption{$\tau = 250$~мкс}
    \end{subfigure}
    \caption{Спектры последовательности прямоугольных импульсов при
             $\nu_{\text{повт}} = 2$~кГц и различных длительностях импульса}
    \label{fig:rect_spectra_tau}
\end{figure}

\clearpage
\subsection{А. Амплитуды и частоты гармоник}
При $T = 1$~мс и $\tau = T/20 = 50$~мкс измерены частоты и амплитуды
первых восьми гармоник; амплитуда первой гармоники $a_1 = \aOne$~мВ.
Теоретические значения рассчитаны по \cref{eq:rect_spectrum,eq:rect_ratio},
результаты приведены в \cref{tab:harmonics} и на \cref{fig:harmonics}.

\begin{table}[!ht]
    \centering
    \caption{Частоты и относительные амплитуды гармоник прямоугольных
             импульсов ($T = 1$~мс, $\tau = 50$~мкс); $\varepsilon$ ---
             относительное отклонение эксперимента от теории}
    \label{tab:harmonics}
    \csvreader[
        separator = semicolon,
        tabularray = {
            colspec = {|c|c|c|c|c|c|c|},
            hlines, vlines,
            row{1} = {font=\bfseries},
        },
        table head = {
            $n$ & $\nu_n^{\text{эксп}}$, кГц & $\nu_n^{\text{теор}}$, кГц
            & $|a_n|$, мВ & $|a_n/a_1|^{\text{эксп}}$
            & $|a_n/a_1|^{\text{теор}}$ & $\varepsilon$, \% \\
        }
    ]%
    {../data/processed/harmonics_processed.csv}%
    {n=\cN, nu_exp_kHz=\cNuE, nu_theor_kHz=\cNuT, a_exp_mV=\cA,
     ratio_exp=\cR, ratio_theor=\cRT, dev_percent=\cDev}%
    {$\cN$ & $\pn{\cNuE}$ & $\pn{\cNuT}$ & $\pn{\cA}$
     & $\pgfmathprintnumber[fixed, zerofill, precision=3]{\cR}$
     & $\pgfmathprintnumber[fixed, zerofill, precision=3]{\cRT}$ & \cDev}
\end{table}

\begin{figure}[!ht]
    \centering
    \begin{tikzpicture}
        \begin{axis}[
            width=0.8\textwidth,
            height=0.5\textwidth,
            xlabel={$\nu$, кГц},
            ylabel={$|a_n/a_1|$},
            xmin=0, xmax=21, ymin=0, ymax=1.05,
            grid=both,
            minor tick num=1,
            legend pos=north east,
            legend cell align=left,
        ]
            \addplot[only marks, mark=*]
            table[col sep=semicolon, x=nu_exp_kHz, y=ratio_exp]
                {../data/processed/harmonics_processed.csv};
            \addlegendentry{Эксперимент}
            % огибающая |sinc(pi nu tau)|, tau = 0,05 мс, нормировка на n = 1
            \addplot[no marks, thick, domain=0.01:21, samples=300]
                {abs(sin(deg(pi*x*0.05))/(pi*x*0.05))/0.995893};
            \addlegendentry{$|\operatorname{sinc}(\pi\nu\tau)|$, $\tau = 50$~мкс}
        \end{axis}
    \end{tikzpicture}
    \caption{Относительные амплитуды гармоник и теоретическая огибающая
             спектра при $T = 1$~мс, $\tau = 50$~мкс}
    \label{fig:harmonics}
\end{figure}

Частоты гармоник совпадают с теоретическими $\nu_n = n/T$ в пределах
точности отсчёта. Относительные амплитуды совпадают с рассчитанными по
\cref{eq:rect_ratio} с точностью лучше $1$~\%, но систематически немного
ниже теоретических: для $n = 2\div8$ отклонение лежит в пределах
от $\hDevMin$ до $\hDevMax$~\%, в среднем $\hDevMean$~\%. Такое систематическое занижение
соответствует чуть более широкому импульсу: аппроксимация огибающей по
\cref{eq:rect_ratio} с $\tau$ в качестве свободного параметра даёт
$\tau = \tauEnv$~мкс, т.\,е. на $\approx1{,}4$~\% больше установленного
значения. Это может быть связано с конечной длительностью фронтов
реального импульса и точностью установки $\tau$ на генераторе.

Первый нуль огибающей при этих параметрах лежит на $1/\tau = 20$~кГц
(\cref{fig:harmonics}); вид спектра в целом соответствует спектрограмме
на \cref{fig:rect_spectra_nu}\,(а), где при выбранном масштабе вершины
линий выходят за границу экрана.

\subsection{А. Спектр меандра в широком диапазоне частот. Наложение частот}
Спектры меандра с частотой $\nu_1 = 50$ и $100$~кГц в полосе
$0\div\screenEdge$~кГц показаны на \cref{fig:meander_screens}. В спектре
присутствуют только нечётные гармоники $\nu_n = n\nu_1$: амплитуды чётных
не превышают $\evenMax$~мВ ($\evenRel\,a_1$), что соответствует
$\tau = T/2$ (\cref{eq:meander}). По положению нечётных гармоник
$\nu_1 = \nuMeandA$~кГц и $\nuMeandB$~кГц. Спектры (а) и (б) сняты при
одинаковых настройках и совпадают.

\begin{figure}[!ht]
    \centering
    \begin{subfigure}{0.32\textwidth}
        \centering
        \labimage{spectrum_am_wide_1.png}{1.0}
        \caption{$\nu_1 = 50$~кГц}
    \end{subfigure}
    \hfill
    \begin{subfigure}{0.32\textwidth}
        \centering
        \labimage{spectrum_am_wide_2.png}{1.0}
        \caption{$\nu_1 = 50$~кГц, повтор}
    \end{subfigure}
    \hfill
    \begin{subfigure}{0.32\textwidth}
        \centering
        \labimage{spectrum_am_wide_3.png}{1.0}
        \caption{$\nu_1 = 100$~кГц}
    \end{subfigure}
    \caption{Спектры меандра в диапазоне $0\div780$~кГц}
    \label{fig:meander_screens}
\end{figure}

\begin{table}[!ht]
    \centering
    \small
    \caption{Нечётные гармоники меандра (частоты исправлены на смещение
             шкалы $\approx 0{,}4$~кГц, найденное по МНК)}
    \label{tab:meander}
    \csvreader[
        separator = semicolon,
        tabularray = {
            colspec = {|c|c|c|c|c|c|c|},
            hlines, vlines,
            row{1} = {font=\bfseries},
        },
        table head = {
            $\nu_1$, кГц & $n$ & $\nu_n$, кГц & $|a_n|$, мВ
            & $|a_n/a_1|$ & $1/n$ & $n|a_n/a_1|$ \\
        }
    ]%
    {../data/processed/meander_processed.csv}%
    {nu1_kHz=\cNuOne, n=\cN, nu_kHz=\cNu, a_mV=\cA, ratio=\cR, inv_n=\cIn,
     n_ratio=\cNR}%
    {$\cNuOne$ & $\cN$ & $\pn{\cNu}$ & $\pn{\cA}$ & $\pn{\cR}$ & $\pn{\cIn}$ & $\pn{\cNR}$}
\end{table}

\begin{figure}[!ht]
    \centering
    \begin{tikzpicture}
        \begin{axis}[
            width=0.8\textwidth,
            height=0.5\textwidth,
            xlabel={$1/n$},
            ylabel={$|a_n/a_1|$},
            xmin=0, xmax=1.05, ymin=0, ymax=1.05,
            grid=both,
            minor tick num=1,
            legend pos=north west,
            legend cell align=left,
        ]
            \addplot[only marks, mark=*]
            table[col sep=semicolon, x=inv_n, y=ratio, restrict expr to domain={\thisrow{nu1_kHz}}{50:50}]
                {../data/processed/meander_processed.csv};
            \addlegendentry{$\nu_1 = 50$~кГц}
            \addplot[only marks, mark=square, mark size=3pt]
            table[col sep=semicolon, x=inv_n, y=ratio, restrict expr to domain={\thisrow{nu1_kHz}}{100:100}]
                {../data/processed/meander_processed.csv};
            \addlegendentry{$\nu_1 = 100$~кГц}
            \addplot[no marks, thick, domain=0:1.05] {x};
            \addlegendentry{теория: $|a_n/a_1| = 1/n$}
        \end{axis}
    \end{tikzpicture}
    \caption{Относительные амплитуды нечётных гармоник меандра}
    \label{fig:meander}
\end{figure}

Относительные амплитуды первых гармоник согласуются с законом $1/n$
(\cref{fig:meander}): при $n = 3$ и $5$ величина $n|a_n/a_1|$ равна
$\nRatioThreeA$ и $\nRatioFiveA$. С ростом номера амплитуды спадают
несколько быстрее, чем $1/n$: при $n = 15$ ($\nu = 750$~кГц)
$n|a_n/a_1| = \nRatioMinA$. Для $\nu_1 = 50$ и $100$~кГц это отклонение
одинаково как функция номера гармоники (а не частоты), поэтому оно,
по-видимому, связано с формой импульсов генератора (конечная длительность
фронтов, масштабируемая вместе с периодом), а не с частотной
характеристикой входного тракта осциллографа.

Кроме гармоник, в спектре есть линии, которых у меандра быть не должно
(\cref{tab:alias}): при $\nu_1 = 50$~кГц --- на частотах $\approx
212{,}\,312,\,\dots,\,712$~кГц, при $\nu_1 = 100$~кГц --- на
$\approx 63,\,138,\,263,\,\dots$~кГц. Все они объясняются наложением
(\cref{eq:alias}) нечётных гармоник, лежащих выше
$\nu_{\text{дискр}}/2$. Экран спектроанализатора охватывает полосу
$0\div\nu_{\text{дискр}}/2$, т.\,е. $\nu_{\text{дискр}}\approx
2\cdot\screenEdge$~кГц; подбор номера гармоники $n$ и кратности $k$ для
каждой линии и последующее усреднение дают
\begin{equation}
    \nu_{\text{дискр}} = \fsamp~\text{кГц},
    \qquad
    \nu_{\text{дискр}}/2 = \fnyq~\text{кГц}.
\end{equation}
Расчётные положения линий наложения совпадают с наблюдаемыми в пределах
погрешности отсчёта ($0{,}75$~кГц), а их амплитуды $n|a_n/a_1|$ плавно
продолжают зависимость из \cref{tab:meander}: это действительно гармоники
с номерами $n = 9\div27$, «перенесённые» в видимую полосу. Одно и то же
значение $\nu_{\text{дискр}}$ описывает оба спектра. Этот пример
иллюстрирует требование $\nu<\nu_{\text{дискр}}/2$: для сигнала с
медленно спадающим спектром часть линий на экране не соответствует
реальным частотам сигнала.

\begin{table}[!ht]
    \centering
    \small
    \caption{Линии наложения в спектрах меандра: наблюдаемая частота
             $\nu'$, номер гармоники-источника $n$, её частота $n\nu_1$,
             кратность $k$ и расчётная частота
             $|n\nu_1 - k\nu_{\text{дискр}}|$}
    \label{tab:alias}
    \csvreader[
        separator = semicolon,
        tabularray = {
            colspec = {|c|c|c|c|c|c|c|c|},
            hlines, vlines,
            row{1} = {font=\bfseries},
        },
        table head = {
            $\nu_1$, кГц & $\nu'$, кГц & $n$ & $n\nu_1$, кГц & $k$
            & $\nu'_{\text{расч}}$, кГц & $a$, мВ & $n a/a_1$ \\
        }
    ]%
    {../data/processed/alias_processed.csv}%
    {nu1_kHz=\cNuOne, nu_obs_kHz=\cObs, n=\cN, nu_n_kHz=\cNn, k=\cK,
     nu_pred_kHz=\cPred, a_mV=\cA, n_ratio=\cNR}%
    {$\cNuOne$ & $\pn{\cObs}$ & $\cN$ & $\pn{\cNn}$ & $\cK$ & $\pn{\cPred}$ & $\pn{\cA}$ & $\pn{\cNR}$}
\end{table}

\clearpage
\subsection{А. Проверка соотношения $\Delta\nu\cdot\tau\simeq1$}
При фиксированном периоде $T = 1$~мс длительность импульса менялась в
диапазоне от $T/50$ до $T/5$; измерялась ширина спектра $\Delta\nu$ ---
от $\nu = 0$ до гармоники с нулевой амплитудой (\cref{tab:width}). Поскольку
гармоники расположены через $\delta\nu = 1$~кГц, погрешность определения
нуля огибающей принята равной $0{,}5$~кГц.

\begin{table}[!ht]
    \centering
    \caption{Ширина спектра при различных длительностях импульса
             ($T = 1$~мс)}
    \label{tab:width}
    \csvreader[
        separator = semicolon,
        tabularray = {
            colspec = {|c|c|c|c|},
            hlines, vlines,
            row{1} = {font=\bfseries},
        },
        table head = {
            $\tau$, мкс & $1/\tau$, кГц & $\Delta\nu$, кГц
            & $\Delta\nu\cdot\tau$ \\
        }
    ]%
    {../data/processed/spectrum_width_processed.csv}%
    {tau_us=\cTau, inv_tau_kHz=\cIt, Dnu_kHz=\cD, Dnu_tau=\cP}%
    {$\cTau$ & $\cIt$ & $\cD$ & $\cP$}
\end{table}

\begin{figure}[!ht]
    \centering
    \begin{tikzpicture}
        \begin{axis}[
            width=0.8\textwidth,
            height=0.5\textwidth,
            xlabel={$1/\tau$, кГц},
            ylabel={$\Delta\nu$, кГц},
            xmin=0, xmax=55, ymin=0, ymax=55,
            grid=both,
            minor tick num=1,
            legend pos=north west,
            legend cell align=left,
        ]
            \addplot[only marks, mark=*,
                     error bars/.cd, y dir=both, y explicit]
            table[col sep=semicolon, x=x, y=y, y error=ey]
                {../data/processed/spectrum_width_processed.csv};
            \addlegendentry{Эксперимент}
            \addplot[no marks, thick, domain=0:55] {\kDnuNum*x};
            \addlegendentry{МНК: $\Delta\nu = k/\tau$, $k = \kDnuVal$}
        \end{axis}
    \end{tikzpicture}
    \caption{Зависимость ширины спектра от обратной длительности импульса}
    \label{fig:width}
\end{figure}

Взвешенный МНК для прямой через начало координат даёт
\begin{equation}
    k_{\Delta\nu} = \kDnu;
\end{equation}
при аппроксимации $\Delta\nu = k/\tau + b$ получаем $k = \kDnuFree$,
$b = \bDnu$~кГц, т.\,е. свободный член в пределах погрешности равен нулю.
Наклон совпадает с теоретическим значением $k = 1$
(\cref{eq:uncertainty_rect}), а произведение $\Delta\nu\cdot\tau$ для всех
точек лежит в пределах $\prodMin\div\prodMax$; отклонения от единицы
(при $\tau = 110\div170$~мкс) объясняются тем, что нуль огибающей
$1/\tau$ попадает между гармониками, и его положение определяется лишь с
точностью до шага $\delta\nu$. Соотношение неопределённостей
$\Delta\nu\cdot\tau \simeq 1$ выполняется в диапазоне длительностей,
различающихся в $10$ раз. Тот же результат даёт и спектр меандра
(\cref{fig:meander_screens}): при $\tau = T/2$ первый нуль огибающей
$1/\tau = 2\nu_1$ совпадает со второй гармоникой, и все чётные гармоники
исчезают.

\clearpage
\subsection{А. Проверка соотношения $\delta\nu\cdot T = 1$}
При фиксированной длительности импульса $\tau = 100$~мкс период повторения
менялся от $2\tau$ до $50\tau$; измерялось расстояние между соседними
гармониками (\cref{tab:spacing}).

\begin{table}[!ht]
    \centering
    \caption{Расстояние между соседними гармониками при различных периодах
             повторения ($\tau = 100$~мкс)}
    \label{tab:spacing}
    \csvreader[
        separator = semicolon,
        tabularray = {
            colspec = {|c|c|c|c|c|c|},
            hlines, vlines,
            row{1} = {font=\bfseries},
        },
        table head = {
            $T$, мс & $T/\tau$ & $1/T$, кГц & $\delta\nu$, кГц
            & $\delta\nu\cdot T$ & примечание \\
        }
    ]%
    {../data/processed/harmonics_spacing_processed.csv}%
    {T_ms=\cT, T_tau=\cTt, inv_T_kHz=\cIT, dnu_kHz=\cD, dnu_T=\cP,
     note=\cNote}%
    {$\cT$ & $\cTt$ & $\cIT$ & $\cD$ & $\cP$ & \cNote}
\end{table}

\begin{figure}[!ht]
    \centering
    \begin{tikzpicture}
        \begin{axis}[
            width=0.8\textwidth,
            height=0.5\textwidth,
            xlabel={$1/T$, кГц},
            ylabel={$\delta\nu$, кГц},
            xmin=0, xmax=1.1, ymin=0, ymax=1.1,
            grid=both,
            minor tick num=1,
            legend pos=north west,
            legend cell align=left,
        ]
            \addplot[only marks, mark=*]
            table[col sep=semicolon, x=x, y=y,
                  restrict expr to domain={\thisrow{use}}{1:1}]
                {../data/processed/harmonics_spacing_processed.csv};
            \addlegendentry{Эксперимент}
            \addplot[no marks, thick, domain=0:1.1] {\kdnuNum*x};
            \addlegendentry{МНК: $\delta\nu = k/T$, $k = \kdnuVal$}
        \end{axis}
    \end{tikzpicture}
    \caption{Зависимость расстояния между гармониками от обратного периода
             повторения (без точки $T = 2\tau$, см. текст)}
    \label{fig:spacing}
\end{figure}

Точка $T = 2\tau$ выпадает: измеренное расстояние вдвое больше $1/T$. Это
не ошибка измерения: при $T = 2\tau$ сигнал является меандром, все чётные
гармоники обращаются в ноль (\cref{eq:meander}), и соседние
\emph{видимые} линии отстоят друг от друга на $2/T = 10$~кГц. Поэтому эта
точка в аппроксимацию не включалась. По остальным шести точкам
\begin{equation}
    k_{\delta\nu} = \kdnu,
\end{equation}
а аппроксимация $\delta\nu = k/T + b$ даёт $k = \kdnuFree$,
$b = \bdnu$~кГц. Наклон совпадает с теоретическим значением $k = 1$,
произведение $\delta\nu\cdot T$ равно единице с точностью $0{,}5$~\%.
Соотношение $\delta\nu = 1/T$ --- точное и от формы сигнала не зависит
(далее оно подтверждается и для цугов).

\clearpage
\subsection{Б. Спектр последовательности цугов}
Генератор настроен на подачу цугов синусоиды с несущей $\nu_0 = 50$~кГц,
частотой повторения $\nu_{\text{повт}} = 1$~кГц ($T = 1$~мс) и числом
периодов в цуге $N = 5$ (длительность цуга $\tau = N/\nu_0 = 100$~мкс).
Затем поочерёдно менялись число периодов $N$, частота несущей $\nu_0$ и
период повторения $T$. В каждом случае измерялись расстояние между
гармониками $\delta\nu$ и ширина спектра $\Delta\nu$ --- от центра спектра
до первого нуля огибающей (\cref{tab:bursts}). Спектрограммы цугов при
различных параметрах приведены на \cref{fig:burst_spectra}.

\begin{figure}[!ht]
    \centering
    \begin{subfigure}{0.48\textwidth}
        \centering
        \labimage{spectrum_burst_T1_N5.png}{1.0}
        \caption{$T = 1$~мс, $\nu_0 = 50$~кГц, $N = 5$}
    \end{subfigure}
    \hfill
    \begin{subfigure}{0.48\textwidth}
        \centering
        \labimage{spectrum_burst_T2_N5.png}{1.0}
        \caption{$T = 2$~мс, $\nu_0 = 50$~кГц, $N = 5$}
    \end{subfigure}

    \vspace{4mm}
    \begin{subfigure}{0.48\textwidth}
        \centering
        \labimage{spectrum_burst_T1_N10.png}{1.0}
        \caption{$T = 1$~мс, $\nu_0 = 50$~кГц, $N = 10$}
    \end{subfigure}
    \hfill
    \begin{subfigure}{0.48\textwidth}
        \centering
        \labimage{spectrum_burst_nu100_N10.png}{1.0}
        \caption{$T = 1$~мс, $\nu_0 = 100$~кГц, $N = 5$}
    \end{subfigure}
    \caption{Спектры последовательности цугов при различных параметрах
             (диапазон частот $23\div83$~кГц). При $\nu_0 = 100$~кГц центр
             спектра лежит за пределами показанного диапазона, видны лишь
             боковые лепестки огибающей}
    \label{fig:burst_spectra}
\end{figure}

\begin{table}[!ht]
    \centering
    \caption{Параметры спектров цугов; в первом столбце указан параметр,
             который менялся в серии. Погрешность $\Delta\nu$ ---
             половина шага гармоник}
    \label{tab:bursts}
    \csvreader[
        separator = semicolon,
        tabularray = {
            colspec = {|c|c|c|c|c|c|c|c|c|},
            hlines, vlines,
            row{1} = {font=\bfseries},
        },
        table head = {
            серия & $\nu_0$, кГц & $T$, мс & $N$ & $\tau$, мкс
            & $\delta\nu$, кГц & $\delta\nu\cdot T$ & $\Delta\nu$, кГц
            & $\Delta\nu\cdot\tau$ \\
        }
    ]%
    {../data/processed/bursts_processed.csv}%
    {series=\cS, nu0_kHz=\cNz, T_ms=\cT, N=\cN, tau_us=\cTau, dnu_kHz=\cd,
     dnu_T=\cdT, Dnu_kHz=\cD, Dnu_tau=\cP}%
    {\cS & $\cNz$ & $\cT$ & $\cN$ & $\cTau$ & $\cd$ & $\cdT$ & $\cD$ & $\cP$}
\end{table}

\begin{figure}[!ht]
    \centering
    \begin{tikzpicture}
        \begin{axis}[
            width=0.8\textwidth,
            height=0.5\textwidth,
            xlabel={$1/\tau = \nu_0/N$, кГц},
            ylabel={$\Delta\nu$, кГц},
            xmin=0, xmax=11, ymin=0, ymax=11,
            grid=both,
            minor tick num=1,
            legend pos=north west,
            legend cell align=left,
        ]
            \addplot[only marks, mark=*,
                     error bars/.cd, y dir=both, y explicit]
            table[col sep=semicolon, x=x, y=y, y error=ey]
                {../data/processed/bursts_width_points.csv};
            \addlegendentry{Эксперимент}
            \addplot[no marks, thick, domain=0:11] {\kBurstNum*x};
            \addlegendentry{МНК: $\Delta\nu = k/\tau$, $k = \kBurstVal$}
        \end{axis}
    \end{tikzpicture}
    \caption{Ширина спектра цугов в зависимости от обратной длительности
             цуга $1/\tau = \nu_0/N$}
    \label{fig:burst_width}
\end{figure}

Результаты серий:
\begin{itemize}
    \item \emph{изменение $N$} ($\nu_0 = 50$~кГц, $T = 1$~мс): спектр
          остаётся симметричным относительно одной и той же частоты
          $\nu_0$, но «сжимается» к ней при увеличении $N$:
          $\Delta\nu = 10;\ 5;\ 3;\ 2{,}5$~кГц при $N = 5;\ 10;\ 15;\ 20$,
          т.\,е. $\Delta\nu = \nu_0/N$ (при $N = 15$ значение
          $\nu_0/N = 3{,}3$~кГц лежит между гармониками, и нуль огибающей
          определяется с точностью до шага $\delta\nu$). Шаг гармоник $\delta\nu = 1$~кГц не
          меняется. На \cref{fig:burst_spectra}\,(а, в) видно, что при
          удвоении $N$ центральный лепесток вдвое сужается, а амплитуда
          центральных линий вдвое растёт;
    \item \emph{изменение $\nu_0$} ($N = 5$, $T = 1$~мс): при
          $\nu_0 = 30;\ 40;\ 50$~кГц ширина спектра $\Delta\nu = 6;\ 8;\
          10$~кГц растёт пропорционально $\nu_0$ (поскольку
          $\tau = N/\nu_0$ уменьшается), а центр спектра смещается вслед за
          частотой несущей. На \cref{fig:burst_spectra}\,(г) при
          $\nu_0 = 100$~кГц нули огибающей видны на $\approx40$, $60$ и
          $80$~кГц, т.\,е. через $\nu_0/N = 20$~кГц;
    \item \emph{изменение $T$} ($\nu_0 = 50$~кГц, $N = 5$): при
          $T = 1;\ 2;\ 4$~мс расстояние между гармониками
          $\delta\nu = 1;\ 0{,}5;\ 0{,}25$~кГц $= 1/T$, а огибающая
          ($\Delta\nu = 10$~кГц) не меняется
          (\cref{fig:burst_spectra}\,(а, б)); при увеличении $T$ амплитуды
          гармоник уменьшаются пропорционально $\tau/T$.
\end{itemize}

Аппроксимация зависимости $\Delta\nu(1/\tau)$ (\cref{fig:burst_width})
прямой через начало координат даёт
\begin{equation}
    k = \kBurst,
\end{equation}
во всех измерениях $\Delta\nu\cdot\tau = 1$ и $\delta\nu\cdot T = 1$ в
пределах точности отсчёта, что подтверждает соотношения
неопределённостей \cref{eq:burst_width} и $\delta\nu = 1/T$. То, что центр
спектра совпадает с частотой несущей, а форма огибающей и её ширина
определяются только длительностью цуга, подтверждает теорему смещения
\cref{eq:shift}: спектр цугов --- это спектр последовательности
прямоугольных импульсов длительности $\tau = N/\nu_0$, перенесённый из
окрестности нуля в окрестность $\nu_0$.

\clearpage
\subsection{Г. Спектр амплитудно-модулированного сигнала}
Генератор настроен на АМ-сигнал с несущей $\nu_0 = 50$~кГц и частотой
модуляции $\nu_{\text{мод}} = 2$~кГц. В спектре (\cref{fig:am_spectrum})
наблюдаются три линии: центральная на частоте $\nu_0$ и две боковые
одинаковой амплитуды на частотах $\nu_0\pm\nu_{\text{мод}}$
(\cref{tab:am}). Обратите внимание: шкала амплитуд на
\cref{fig:am_spectrum} начинается не с нуля, а с $19{,}17$~мВ.

\begin{figure}[!ht]
    \centering
    \labimage{spectrum_am.png}{0.7}
    \caption{Спектр АМ-сигнала: $\nu_0 = 50$~кГц,
             $\nu_{\text{мод}} = 2$~кГц}
    \label{fig:am_spectrum}
\end{figure}

\begin{table}[!ht]
    \centering
    \caption{Спектральные линии АМ-сигнала}
    \label{tab:am}
    \csvreader[
        separator = semicolon,
        tabularray = {
            colspec = {|c|c|c|c|},
            hlines, vlines,
            row{1} = {font=\bfseries},
        },
        table head = {
            линия & $\nu$, кГц & $a$, мВ & $a/a_{\text{осн}}$ \\
        }
    ]%
    {../data/processed/am_processed.csv}%
    {line=\cL, nu_kHz=\cNu, a_mV=\cA, ratio=\cR}%
    {\cL & $\cNu$ & $\cA$ & $\cR$}
\end{table}

По положению линий $\nu_0 = \amNuzero$~кГц,
$\nu_{\text{мод}} = \amNumod$~кГц, что совпадает с установленными
значениями. Отношение амплитуд
\begin{equation}
    \frac{a_{\text{бок}}}{a_{\text{осн}}} = \amRatio
    \quad\Rightarrow\quad
    m = \frac{2a_{\text{бок}}}{a_{\text{осн}}} = \amM,
\end{equation}
т.\,е. спектр снят при глубине модуляции $m = 100$~\%. Боковые линии
симметричны: их амплитуды совпадают в пределах цены деления
($1{,}5$~мВ). При $m = 1$ соотношение \cref{eq:am_ratio} предсказывает
$a_{\text{бок}}/a_{\text{осн}} = 0{,}5$; измеренное значение отличается
от него на $0{,}6$~\%, т.\,е. на $1{,}8$ погрешности, что соответствует
точности установки $m$ на генераторе.

\clearpage
\subsection{Е. Фильтрация сигнала $RC$-цепочкой}
Была собрана схема фильтра низких частот (\cref{fig:rc_scheme}) и на её вход
подана последовательность прямоугольных импульсов. Получены осциллограмма
сигнала на выходе цепочки и его спектр (\cref{fig:filter}). Прямое
измерение коэффициентов фильтрации $K_n$ (со спектром входного сигнала,
снятым отдельно) провести не удалось; постоянная времени определена по
форме выходного сигнала.

\begin{figure}[!ht]
    \centering
    \begin{subfigure}{0.48\textwidth}
        \centering
        \labimage{filter_signal.png}{1.0}
        \caption{осциллограмма (вход «Перем.», мкс)}
    \end{subfigure}
    \hfill
    \begin{subfigure}{0.48\textwidth}
        \centering
        \labimage{filter_spectrum.png}{1.0}
        \caption{спектр (МГц)}
    \end{subfigure}
    \caption{Сигнал на выходе $RC$-цепочки и его спектр}
    \label{fig:filter}
\end{figure}

\paragraph{Осциллограмма.}
Сигнал на конденсаторе имеет вид чередующихся коротких участков нарастания
и длинных участков почти линейного спада: за время импульса конденсатор
заряжается, в паузе --- разряжается. По положению вершин период
$T = \filtT$~мкс. Каждый участок спада аппроксимирован экспонентой
\cref{eq:rc_transient} со свободными $U_\infty$, $U(0)$ и $\tau_{RC}$;
среднее по трём периодам
\begin{equation}
    \tau_{RC} = \filtRC~\text{мкс},
    \qquad
    \nu_{RC} = \frac{1}{2\pi\tau_{RC}} = \filtNuRC~\text{кГц}.
\end{equation}
При известном $\tau_{RC}$ аппроксимация участков нарастания даёт
высоту входных импульсов $U_0 = \filtUzero$~В (генератор выдаёт импульсы
размахом $4$~В) и их длительность $\tau = \filtTau$~мкс $\approx T/20$ ---
именно такие параметры рекомендует методика. Поскольку
$\tau_{RC}\approx 1{,}2\,T \gg \tau$, выходной сигнал имеет малую (около
$4$~\% от $U_0$) амплитуду, а спады почти линейны.

\paragraph{Спектр.}
Спектр выходного сигнала дискретен: линии стоят на частотах, кратных
$\nu_1 = \filtNuOne$~кГц ($T = \filtTspec$~мкс, в согласии с
осциллограммой). Первая гармоника лежит далеко за частотой среза фильтра:
$2\pi\nu_1\tau_{RC} = \filtWnu \gg 1$, т.\,е. цепочка работает в режиме
интегрирования. Поэтому, хотя у входных импульсов с $\tau\approx T/20$
амплитуды первых десяти гармоник почти одинаковы, на выходе они убывают
примерно как $1/n$.

Для сравнения с \cref{eq:rc_filter} относительные коэффициенты
фильтрации вычислены как
\begin{equation}
    \frac{K_n}{K_1} = \frac{a_n/a_1}{\bigl(a_n/a_1\bigr)_{\text{вх}}},
\end{equation}
где $(a_n/a_1)_{\text{вх}}$ рассчитано по \cref{eq:rect_ratio} для
найденного $\tau/T = \filtQ$ (\cref{tab:filter}, \cref{fig:filter_K}).

\begin{table}[!ht]
    \centering
    \caption{Гармоники сигнала на выходе $RC$-цепочки и относительный
             коэффициент фильтрации ($\tau_{RC} = \filtRC$~мкс)}
    \label{tab:filter}
    \csvreader[
        separator = semicolon,
        tabularray = {
            colspec = {|c|c|c|c|c|c|c|},
            hlines, vlines,
            row{1} = {font=\bfseries},
        },
        table head = {
            $n$ & $\nu_n$, кГц & $a_n$, мВ & $(a_n/a_1)_{\text{вх}}$
            & $(K_n/K_1)_{\text{эксп}}$ & $(K_n/K_1)_{\text{теор}}$
            & $\varepsilon$, \% \\
        }
    ]%
    {../data/processed/filter_processed.csv}%
    {n=\cN, nu_kHz=\cNu, a_mV=\cA, ain_rel=\cAin, K_exp=\cKe,
     K_theor=\cKt, dev_percent=\cDev}%
    {$\cN$ & $\pn{\cNu}$ & $\pn{\cA}$ & $\pn{\cAin}$ & $\pn{\cKe}$ & $\pn{\cKt}$ & \cDev}
\end{table}

\begin{figure}[!ht]
    \centering
    \begin{tikzpicture}
        \begin{axis}[
            width=0.8\textwidth,
            height=0.5\textwidth,
            xlabel={$\nu$, кГц},
            ylabel={$K_n/K_1$},
            xmin=0, xmax=3600, ymin=0, ymax=1.1,
            grid=both,
            minor tick num=1,
            legend pos=north east,
            legend cell align=left,
            scaled x ticks=false,
            xticklabel style={/pgf/number format/1000 sep={}},
        ]
            \addplot[only marks, mark=*]
            table[col sep=semicolon, x=nu_kHz, y=K_exp]
                {../data/processed/filter_processed.csv};
            \addlegendentry{Эксперимент}
            \addplot[no marks, thick, domain=330:3600, samples=200]
                {sqrt(1+(2*pi*\filtNuOneNum*\filtRCNum/1000)^2)
                 /sqrt(1+(2*pi*x*\filtRCNum/1000)^2)};
            \addlegendentry{\cref{eq:rc_filter}, $\tau_{RC} = \filtRCVal$~мкс}
        \end{axis}
    \end{tikzpicture}
    \caption{Относительный коэффициент фильтрации $RC$-цепочки}
    \label{fig:filter_K}
\end{figure}

Экспериментальные точки в целом следуют теоретической зависимости
$K\propto1/\nu$: для $n = 2,\,3,\,6,\,9$ отклонение не превышает
$6$~\%, однако для отдельных гармоник ($n = 4,\,5,\,7,\,8,\,10$) оно
достигает $\filtDevMax$~\% (медианное отклонение $\filtDevMedian$~\%).
Причины --- спектр входного сигнала не измерялся, а рассчитывался для
идеальных импульсов (форма реальных коротких импульсов может заметно отличаться от
прямоугольной), и погрешность БПФ при коротком окне для линий с амплитудой в
несколько мВ. Кроме того, при $2\pi\nu\tau_{RC}\gg1$ отношение
$K_n/K_1\approx1/n$ практически не зависит от $\tau_{RC}$, поэтому
определять $\tau_{RC}$ по этим данным нельзя; её значение надёжно
определено по осциллограмме. Качественно спектр подтверждает, что
$RC$-цепочка работает как фильтр низких частот, подавляющий гармоники
пропорционально их частоте.

\section{Обсуждение результатов}
\begin{itemize}
    \item Частоты гармоник последовательности прямоугольных импульсов
          совпадают с теоретическими $\nu_n = n/T$. Относительные амплитуды
          первых восьми гармоник совпадают с \cref{eq:rect_ratio} с
          точностью лучше $1$~\%; небольшое систематическое занижение
          (в среднем $\hDevMean$~\%) соответствует эффективной длительности
          импульса $\tau = \tauEnv$~мкс вместо установленных $50$~мкс.
    \item Соотношения неопределённостей подтверждены: наклоны прямых
          $\Delta\nu(1/\tau)$ и $\delta\nu(1/T)$ равны $\kDnu$ и $\kdnu$
          при теоретическом значении 1. Погрешность $\Delta\nu$ определяется
          дискретностью спектра: нуль огибающей можно найти лишь с
          точностью до шага гармоник. Точка $T = 2\tau$ в зависимости
          $\delta\nu(1/T)$ выпадает закономерно: у меандра нет чётных
          гармоник, и видимые линии отстоят на $2/T$.
    \item В спектре меандра в полосе $0\div780$~кГц наряду с нечётными
          гармониками ($|a_n/a_1|\approx1/n$) наблюдались линии наложения
          гармоник с $n = 9\div27$; по ним определена частота дискретизации
          спектроанализатора $\nu_{\text{дискр}} = \fsamp$~кГц. Это наглядно
          показывает, почему спектр можно считать достоверным лишь при
          $\nu<\nu_{\text{дискр}}/2$.
    \item Теорема смещения подтверждена: спектр цугов симметричен
          относительно частоты несущей и повторяет спектр прямоугольных
          импульсов длительности $\tau = N/\nu_0$; $\Delta\nu = \nu_0/N$
          ($k = \kBurst$) и $\delta\nu = 1/T$ во всех сериях.
    \item Для АМ-сигнала положения линий $\nu_0$ и
          $\nu_0\pm\nu_{\text{мод}}$ совпадают с заданными, а
          $a_{\text{бок}}/a_{\text{осн}} = \amRatio$ соответствует
          $m = \amM$. Сохранён спектр только при одной глубине модуляции,
          поэтому зависимость $a_{\text{бок}}/a_{\text{осн}}(m)$ не
          построена.
    \item Для $RC$-цепочки по осциллограмме определена
          $\tau_{RC} = \filtRC$~мкс; спектр выходного сигнала согласуется
          с режимом интегрирования ($K\propto1/\nu$), но из-за отсутствия
          измеренного спектра входного сигнала расхождение для отдельных
          гармоник достигает $\filtDevMax$~\%.
    \item Ограничения метода: спектр вычисляется только по видимой на экране
          части сигнала, что вносит «краевые эффекты» (слабые побочные пики
          на спектрограммах части А) и ограничивает разрешение по частоте;
          дискретизация ограничивает доступный диапазон частот условием
          $\nu < \nu_{\text{дискр}}/2$.
\end{itemize}

\section{Заключение}
В работе изучены спектры периодических сигналов различной формы и проверены
основные соотношения фурье-анализа:
\begin{itemize}
    \item при увеличении частоты повторения расстояние между компонентами
          спектра увеличивается, при увеличении длительности импульса
          ширина спектра уменьшается;
    \item частоты гармоник последовательности прямоугольных импульсов
          совпадают с $n/T$, относительные амплитуды --- с
          \cref{eq:rect_ratio} с точностью лучше $1$~\%; при $\tau = T/2$
          чётные гармоники исчезают, нечётные убывают как $1/n$;
    \item проверены соотношения неопределённостей $\Delta\nu = k/\tau$ и
          $\delta\nu = k/T$: получены наклоны $\kDnu$ и $\kdnu$ при
          теоретическом значении $k = 1$;
    \item на примере спектра последовательности цугов подтверждена теорема
          смещения: центр спектра совпадает с частотой несущей $\nu_0$,
          шаг гармоник равен $1/T$, а для ширины спектра
          $\Delta\nu = k\nu_0/N$ получено $k = \kBurst$;
    \item для амплитудно-модулированного сигнала при $m = 1$ получено
          $a_{\text{бок}}/a_{\text{осн}} = \amRatio$ (теория $0{,}5$);
    \item обнаружено наложение частот в спектре меандра и определена
          частота дискретизации спектроанализатора
          $\nu_{\text{дискр}} = \fsamp$~кГц;
    \item определена постоянная времени $RC$-цепочки
          $\tau_{RC} = \filtRC$~мкс и показано, что она работает как фильтр
          низких частот.
\end{itemize}

\end{document}
